\begin{abstract}
\noindent
We describe a method that decides, for every flight and every instant, whether a paraglider
pilot flies in a group or alone and, in the first case, with which other pilots. Two pilots are
in contact when one of them is where the other was, or will be, a few minutes earlier or later.
This is the proximity with a local delay of Li, Wu and Crofoot~\cite{li2013}, adapted to thermal
flight in three ways: altitude acts as a clock in a thermal and keeps the delay defined where
the path turns on itself; headings are compared only in glide; the maximum delay is tied to the
visual range. At each instant the groups are the connected components of the contact graph. As
in the grouping structure of Buchin et al.~\cite{buchin2015}, a configuration is a fixed set of
flights over a maximal interval, and an episode is a chain of configurations that share members.
Contacts are found in two stages, a coarse filter that provably loses none and an exact scan;
their minimum duration is calibrated on a null model that shifts tracks in time. The output is a
few tables that label every instant of every flight, so that group flight can be compared with
solo flight.
\end{abstract}
